% !TEX root = ../main.tex
\chapter{Hasil Percobaan dan Analisis}

\section{Hasil Preprocessing}

Tahap \textit{preprocessing} mengubah 256 baris \textit{raw data} menjadi \textit{event log} yang siap dianalisis. Dari tahap ini teridentifikasi 52 \textit{case} pemeliharaan, empat di antaranya terbentuk melalui pembangkitan \textit{Case ID} otomatis yang mencakup 21 baris pekerjaan. \textit{Compound activity splitting} menaikkan jumlah \textit{event} dari 255 pada \textit{event log as-is} menjadi 308 pada \textit{event log enriched}, sementara imputasi \textit{timestamp} menyentuh 110 \textit{event}. Ringkasan tahap ini disajikan pada Tabel~\ref{tab:ringkasan-preprocessing}.

\begin{table}[H]
\centering
\caption{Ringkasan Hasil Preprocessing}
\label{tab:ringkasan-preprocessing}
\begin{tabular}{lr}
\toprule
\textbf{Indikator} & \textbf{Nilai} \\
\midrule
Baris \textit{raw data} & 256 \\
Jumlah \textit{cases} & 52 \\
\textit{Cases} AUTO GENERATED       & 4 (21 baris) \\
\textit{Event} pada \textit{event log as-is} & 255 \\
\textit{Event} pada \textit{event log enriched} & 308 \\
\textit{Timestamp} yang diimputasi & 110 \\
\textit{Compound activity splitting} & 99 \\
\bottomrule
\end{tabular}
\end{table}

\section{Hasil Anotasi LLM}

Model bahasa memberi anotasi pada seluruh 308 \textit{event} di 52 \textit{case}. Seluruh \textit{event} terdistribusi ke dalam dua belas kategori \texttt{action\_type}. TESTING menempati frekuensi tertinggi dengan 53 \textit{event} (17,21\%), diikuti SAFETY\_INDUCTION dan CLEANING yang masing-masing mencakup 52 \textit{event} (16,88\%), serta INSPECTION dengan 51 \textit{event} (16,56\%). Komposisi ini mencerminkan karakteristik pekerjaan \textit{overhaul} yang didominasi pengujian, pengarahan keselamatan, pembersihan, dan pemeriksaan. Sebaliknya, REPAIR, OTHER, dan LUBRICATION muncul jarang sesuai sifat pekerjaannya. Distribusi lengkap disajikan pada Tabel~\ref{tab:distribusi-action-type}.

\begin{table}[H]
\centering
\caption{Distribusi Action Type pada Event Log Enriched}
\label{tab:distribusi-action-type}
\footnotesize
\setlength{\tabcolsep}{5pt}
\renewcommand{\arraystretch}{1.05}
\begin{tabular}{@{}l r r r@{}}
\toprule
\textbf{Action Type} & \textbf{Jumlah \textit{Event}} & \textbf{Persentase} & \textbf{Rerata \textit{Conf.}} \\
\midrule
TESTING & 53 & 17,21\% & 0,96 \\
SAFETY\_INDUCTION & 52 & 16,88\% & 1,00 \\
CLEANING & 52 & 16,88\% & 0,94 \\
INSPECTION & 51 & 16,56\% & 0,93 \\
MEASUREMENT & 45 & 14,61\% & 0,95 \\
REPLACEMENT & 20 & 6,49\% & 0,95 \\
DISASSEMBLY & 12 & 3,90\% & 0,97 \\
ADJUSTMENT & 9 & 2,92\% & 0,86 \\
ASSEMBLY & 7 & 2,27\% & 0,98 \\
REPAIR & 4 & 1,30\% & 0,93 \\
OTHER & 2 & 0,65\% & 0,97 \\
LUBRICATION & 1 & 0,32\% & 0,90 \\
\midrule
\textbf{Total} & \textbf{308} & \textbf{100\%} & \\
\bottomrule
\end{tabular}
\end{table}

Di luar kategori tindakan, model mengembalikan atribut kontekstual untuk tiap \textit{event}, meliputi nama komponen, subsistem, area WBS, tingkat \textit{confidence}, \textit{semantic basis}, dan penanda hasil \textit{compound activity splitting}. Contoh hasil anotasi ditampilkan pada Tabel~\ref{tab:contoh-anotasi}.

\begin{table}[H]
\centering
\caption{Contoh Hasil Anotasi Aktivitas}
\label{tab:contoh-anotasi}
\scriptsize
\setlength{\tabcolsep}{3pt}
\renewcommand{\arraystretch}{1.1}
\begin{tabular}{
  >{\raggedright\arraybackslash}p{1.45cm}
  >{\raggedright\arraybackslash}p{2.55cm}
  >{\raggedright\arraybackslash}p{2.1cm}
  @{\hspace{14pt}}
  >{\raggedright\arraybackslash}p{1.7cm}
  >{\raggedright\arraybackslash}p{2.1cm}
  @{\hspace{14pt}}
  >{\centering\arraybackslash}p{0.7cm}
}
\toprule
\textbf{Case ID} & \textbf{Aktivitas (\textit{raw})} & \textbf{Action Type} & \textbf{Komponen} & \textbf{Sub-sistem} & \textbf{Conf.} \\
\midrule
AUTO\_001 & Safety Induction & SAFETY\_INDUCTION & GENERAL & GENERAL & 1,00 \\
AUTO\_001 & PEMBERSIHAN & CLEANING & GENERAL & GENERAL & 0,95 \\
AUTO\_001 & PEMERIKSAAN & INSPECTION & GENERAL & GENERAL & 0,95 \\
AUTO\_002 & IR Test & MEASUREMENT & PANEL 52G & ELECTRICAL\_AUX & 1,00 \\
AUTO\_004 & LEPAS THERMOSTAT & DISASSEMBLY & THERMOSTAT & COOLING\_SYSTEM & 1,00 \\
AUTO\_004 & PEMASANGAN THERMOSTAT & ASSEMBLY & THERMOSTAT & COOLING\_SYSTEM & 1,00 \\
\bottomrule
\end{tabular}
\end{table}

\section{Hasil Process Discovery}

Kedua skenario menghasilkan \textit{Petri net} dengan karakteristik yang jauh berbeda, terutama pada jumlah aktivitas unik yang menjadi basis model. Ringkasan ukuran kedua model disajikan pada Tabel~\ref{tab:perbandingan-ukuran-model}.

\begin{table}[H]
\centering
\caption{Perbandingan Ukuran Model As-Is dan Enriched AI}
\label{tab:perbandingan-ukuran-model}
\footnotesize
\setlength{\tabcolsep}{5pt}
\renewcommand{\arraystretch}{1.05}
\begin{tabular}{@{}l r r r r@{}}
\toprule
\textbf{Metrik} & \textbf{As-Is} & \textbf{Enriched AI} & $\boldsymbol{\Delta}$ & \textbf{Perubahan} \\
\midrule
Cases & 52 & 52 & +0 & +0,00\% \\
Events & 255 & 308 & +53 & +20,78\% \\
Activities & 112 & 12 & $-$100 & $-$89,29\% \\
Variants & 36 & 30 & $-$6 & $-$16,67\% \\
Places & 81 & 31 & $-$50 & $-$61,73\% \\
Transitions & 119 & 43 & $-$76 & $-$63,87\% \\
\bottomrule
\end{tabular}
\end{table}

\subsection{Model As-Is}

Model \textit{as-is} dibangun dari 112 label aktivitas unik yang berasal langsung dari \textit{raw activity description}. Banyaknya label menghasilkan struktur padat, yaitu 81 \textit{place}, 119 \textit{transition}, dan 36 varian \textit{trace}. Struktur semacam ini sukar dibaca karena satu tindakan yang sama terwakili oleh banyak label akibat keragaman penulisan teknisi. Gambar~\ref{fig:petri_asis} memperlihatkan model \textit{spaghetti} yang terbentuk.

\begin{figure}[H]
    \centering
    \hspace*{-1.2cm}\makebox[\textwidth][c]{\includegraphics[width=1.5\textwidth]{bab4/petri_net_asis_2.png}}
    \caption{Visualisasi Petri Net Model As-Is}
    \label{fig:petri_asis}
\end{figure}

\subsection{Model Enriched AI}

Model \textit{enriched} dibangun dari dua belas kategori \texttt{action\_type}. Abstraksi ini memangkas aktivitas unik sebesar 89,29\% dan menyusutkan struktur menjadi 31 \textit{place}, 43 \textit{transition}, dan 30 varian \textit{trace}. Alur pemeliharaan menjadi terbaca meskipun jumlah \textit{event} justru bertambah akibat \textit{compound activity splitting}. Gambar~\ref{fig:petri_enriched} menampilkan modelnya.

\begin{figure}[H]
    \centering
    \hspace*{-1cm}\makebox[\textwidth][c]{\includegraphics[width=1.42\textwidth]{bab4/petri_net_enriched_2.png}}
    \caption{Visualisasi Petri Net Model Enriched AI}
    \label{fig:petri_enriched}
\end{figure}

Untuk membaca keterhubungan antar tindakan, Tabel~\ref{tab:dfg-teratas} menyajikan enam relasi \textit{directly-follows} (DFG) berfrekuensi tertinggi beserta rata-rata selang penyelesaiannya. Pola ini menegaskan alur baku pekerjaan, mulai dari pengarahan keselamatan, pengukuran, pembersihan, dan pemeriksaan, hingga pengujian akhir.

\begin{table}[H]
\centering
\caption{DFG (Enriched AI)}
\label{tab:dfg-teratas}
\footnotesize
\setlength{\tabcolsep}{5pt}
\renewcommand{\arraystretch}{1.05}
\begin{tabular}{@{}l l r r@{}}
\toprule
\textbf{Dari} & \textbf{Ke} & \textbf{Frekuensi} & \textbf{Jeda (jam)} \\
\midrule
CLEANING & INSPECTION & 21 & 6,25 \\
MEASUREMENT & CLEANING & 19 & 13,44 \\
SAFETY\_INDUCTION & MEASUREMENT & 18 & 7,04 \\
INSPECTION & TESTING & 16 & 1,47 \\
SAFETY\_INDUCTION & CLEANING & 13 & 18,53 \\
SAFETY\_INDUCTION & INSPECTION & 13 & 8,13 \\
\bottomrule
\end{tabular}
\end{table}

% [Placeholder Gambar 4.4] Directly-Follows Graph atau heatmap frekuensi antar action_type pada model Enriched AI.

\section{Hasil Evaluasi dan Pembahasan}

\subsection{Conformance Checking}

Kualitas kedua model diukur pada empat dimensi \textit{conformance}, yaitu \textit{fitness}, \textit{precision}, \textit{generalization}, dan \textit{simplicity}. Hasil pengukuran disajikan pada Tabel~\ref{tab:perbandingan-conformance}.

\begin{table}[H]
\centering
\caption{Perbandingan Conformance As-Is dan Enriched AI}
\label{tab:perbandingan-conformance}
\footnotesize
\setlength{\tabcolsep}{6pt}
\renewcommand{\arraystretch}{1.05}
\begin{tabular}{@{}l r r r@{}}
\toprule
\textbf{Metrik} & \textbf{As-Is} & \textbf{Enriched AI} & $\boldsymbol{\Delta}$ \\
\midrule
Fitness & 0,9833 & 0,9570 & $-$0,0263 \\
Precision & 0,8716 & 0,5063 & $-$0,3653 \\
Generalization & 0,0807 & 0,7079 & $+$0,6272 \\
Simplicity & 0,7246 & 0,6727 & $-$0,0519 \\
\bottomrule
\end{tabular}
\end{table}

% [Placeholder Gambar 4.5] Radar chart atau diagram batang perbandingan metrik conformance As-Is vs Enriched AI.

Hasil ini memperlihatkan pertukaran nilai yang khas antara kedua model. Model \textit{as-is} unggul pada \textit{fitness} (0,9833) dan \textit{precision} (0,8716), namun keunggulan tersebut wajar karena model dibangun dari 112 \textit{raw label} yang nyaris memetakan setiap \textit{event} satu per satu. Konsekuensinya, model \textit{as-is} \textit{overfitting} terhadap data, tampak dari \textit{generalization} yang sangat rendah, yaitu 0,0807. Model semacam ini menghafal \textit{trace} yang ada tetapi rapuh ketika menghadapi perilaku baru.

Model \textit{enriched} menempuh arah yang berbeda. Abstraksi dua belas kategori menurunkan \textit{precision} menjadi 0,5063 dan sedikit menurunkan \textit{fitness} menjadi 0,9570, tetapi menaikkan \textit{generalization} menjadi 0,7079. Kenaikan ini mengindikasikan model lebih mampu mengakomodasi variasi proses yang wajar tanpa harus melihat setiap kemungkinan \textit{trace}, sedangkan \textit{fitness} yang tetap tinggi menunjukkan model masih me-\textit{replay} hampir seluruh \textit{trace} aktual pemeliharaan.

\subsection{Kesederhanaan dan Keterbacaan Model Proses}

Penilaian kesederhanaan menuntut kehati-hatian karena metrik \textit{simplicity} PM4Py justru turun dari 0,7246 pada model \textit{as-is} menjadi 0,6727 pada model \textit{enriched}. Penurunan tersebut tidak bertentangan dengan klaim keterbacaan. Metrik ini mengukur rata-rata derajat keterhubungan (\textit{arc degree}) tiap simpul pada \textit{Petri net}, sehingga nilainya menurun ketika abstraksi memunculkan percabangan dan paralelisme yang lebih eksplisit di antara dua belas kategori tindakan. Dengan kata lain, metrik ini sensitif terhadap kepadatan koneksi lokal, bukan terhadap ukuran model secara keseluruhan.

Ukuran struktural justru memperlihatkan penyederhanaan yang besar. Aktivitas unik berkurang 89,29\% dari 112 menjadi 12, jumlah \textit{place} menyusut 61,73\% dari 81 menjadi 31, dan jumlah \textit{transition} turun 63,87\% dari 119 menjadi 43. Model \textit{as-is} yang memuat ratusan label mentah menghasilkan struktur \textit{spaghetti} yang sukar ditelusuri analis, sedangkan model \textit{enriched} menampilkan alur pemeliharaan yang ringkas dan terbaca dalam satu pandangan. Keterbacaan bagi manusia lebih tepat diwakili oleh pengurangan ukuran dan keragaman label ini daripada oleh satu skor \textit{simplicity} tunggal.

\subsection{Validasi Output Anotasi LLM}
\label{subsec:validasi-anotasi}

Seluruh 308 \textit{event} berhasil teranotasi oleh model, sehingga dasar kepercayaan terhadap hasil perlu dinyatakan secara eksplisit. Validasi bertumpu pada tiga lapis. Pertama, tiap keluaran menyertakan skor \textit{confidence}, dan sebagian besar \textit{event} berada pada rentang 0,95 hingga 1,00 sebagaimana tampak pada Tabel~\ref{tab:contoh-anotasi}. Kedua, ambang \textit{confidence} 0,45 diterapkan untuk menandai \textit{event} anomali, dan tidak ada \textit{event} yang jatuh di bawahnya. Ketiga, konsistensi keluaran dijaga melalui \textit{temperature} rendah sebesar 0,1 pada model \texttt{gemini-3.1-flash-lite} yang menekan variasi acak.

Skor \textit{confidence} yang tinggi mencerminkan keyakinan model, bukan bukti kebenaran mutlak. Pemeriksaan informal pada Tabel~\ref{tab:contoh-anotasi} memperlihatkan pemetaan yang masuk akal secara domain, misalnya ``IR Test'' menjadi MEASUREMENT dan ``LEPAS THERMOSTAT'' menjadi DISASSEMBLY. Akurasi anotasi belum terukur secara kuantitatif, dan keterbatasan ini diakui secara terbuka, sehingga pemeriksaan menyeluruh oleh ahli pemeliharaan diusulkan sebagai langkah lanjutan.

Kualitas anotasi juga bertumpu pada beberapa perilaku model yang layak dicatat. Model membedakan aktivitas beraroma pengujian secara semantik, misalnya ``IR Test'' dipetakan ke MEASUREMENT karena menghasilkan nilai ukur, sedangkan ``Function Test'' dipetakan ke TESTING karena bersifat verifikasi fungsi. Model juga memecah \textit{compound activity}, yaitu satu baris pekerjaan yang memuat beberapa tindakan, menjadi \textit{event} terpisah dan menyumbang 99 \textit{event} pemecahan. Di samping itu, lima kategori teratas menampung sekitar 82\% \textit{event}, sehingga kesimpulan pada kategori berfrekuensi rendah seperti LUBRICATION perlu kehati-hatian.

\section{Analisis Performa dan Manfaat Operasional}

Model \textit{enriched} membuka pembacaan proses bisnis pemeliharaan yang tidak tampak pada model \textit{as-is}. Alur baku \textit{overhaul} terbentuk dari pengarahan keselamatan, tahap diagnosis (CLEANING, INSPECTION, MEASUREMENT), tahap perbaikan (DISASSEMBLY, REPAIR, REPLACEMENT, ASSEMBLY), hingga TESTING sebagai verifikasi akhir.

Manfaat pertama adalah efisiensi analisis proses. Reduksi aktivitas unik sebesar 89,29\% dari 112 menjadi 12 dan transisi sebesar 63,87\% mengubah model \textit{spaghetti} menjadi alur dua belas tahap yang terbaca.

Manfaat kedua adalah penajaman prioritas sumber daya. CLEANING menyerap total waktu terbesar sehingga efisiensi penyiapan alat dan material berdampak paling luas, sedangkan REPAIR dan ADJUSTMENT paling lama per kejadian sehingga menuntut perencanaan suku cadang dan keahlian khusus.

Manfaat ketiga adalah perbaikan proses bisnis berdasarkan anomali proses yang teridentifikasi, sebagaimana dirinci pada Subbab~\ref{subsec:temuan-anomali}.

\subsection{Identifikasi Bottleneck}
\label{subsec:identifikasi-bottleneck}

Durasi pengerjaan dianalisis pada tingkat kategori tindakan sebagai jam kerja efektif sesuai SOP shift (08.00 sampai 16.30, setara 8,5 jam per hari). Akhir pekan tetap dihitung sebagai hari kerja karena \textit{overhaul} berjalan menerus, sedangkan SAFETY\_INDUCTION dan pengujian pascapemeliharaan ditetapkan sebagai blok tetap 30 menit. Tabel~\ref{tab:durasi-kategori} menyajikan jumlah kejadian, rata-rata, dan total durasi tiap kategori.

\begin{table}[H]
\centering
\caption{Durasi Kerja per Kategori Action Type}
\label{tab:durasi-kategori}
\footnotesize
\setlength{\tabcolsep}{5pt}
\renewcommand{\arraystretch}{1.05}
\begin{tabular}{@{}l r r r@{}}
\toprule
\textbf{Aktivitas} & \textbf{Jumlah \textit{Event}} & \textbf{Rata-rata Durasi Kerja (jam)} & \textbf{Total Durasi Kerja (jam)} \\
\midrule
CLEANING & 52 & 7,23 & 375,96 \\
MEASUREMENT & 45 & 6,27 & 282,15 \\
INSPECTION & 51 & 5,39 & 274,89 \\
ADJUSTMENT & 9 & 21,56 & 194,04 \\
REPLACEMENT & 20 & 7,35 & 147,00 \\
REPAIR & 4 & 32,00 & 128,00 \\
DISASSEMBLY & 12 & 4,83 & 57,96 \\
LUBRICATION & 1 & 42,50 & 42,50 \\
ASSEMBLY & 7 & 4,86 & 34,02 \\
TESTING & 53 & 0,50 & 26,50 \\
SAFETY\_INDUCTION & 52 & 0,50 & 26,00 \\
OTHER & 2 & 10,00 & 20,00 \\
\bottomrule
\end{tabular}
\end{table}

Terdapat dua jenis \textit{bottleneck}. Secara akumulatif, CLEANING menyumbang total waktu terbesar (375,96 jam kerja), diikuti MEASUREMENT dan INSPECTION, karena frekuensinya tinggi. Secara per kejadian, REPAIR dan ADJUSTMENT paling lama (rata-rata 32,00 jam dan 21,56 jam) meski jarang terjadi. Pola ini mengarahkan dua strategi, yaitu efisiensi berbasis volume untuk pembersihan dan pemeriksaan serta perencanaan kesiapan sumber daya untuk perbaikan dan penyetelan.

Interpretasi ini bersifat indikatif. Total durasi kerja seluruh kategori sekitar 1.609 jam, dibanding sekitar 3.071 jam bila memakai waktu kalender, dan sekitar 1.244 jam jika akhir pekan diperlakukan sebagai hari libur. Sebanyak 110 dari 308 \textit{timestamp} merupakan hasil imputasi penjadwalan, dan kategori berfrekuensi sangat kecil seperti LUBRICATION tidak dapat menjadi dasar kesimpulan yang kuat.

\subsection{Temuan Anomali Proses dari Data Pemeliharaan yang Diuji}
\label{subsec:temuan-anomali}

\textit{Process mining} juga mengungkap anomali proses, yaitu penyimpangan urutan aktivitas dari alur ideal. Anomali ini bersumber dari urutan dan pengulangan pada \textit{event log}, berbeda dari anomali pelabelan pada Subbab~\ref{subsec:validasi-anotasi}. Empat pola berikut menjadi peluang perbaikan utama.

Pertama, intervensi tanpa diagnosis. Sebanyak 8 dari 52 \textit{case} melakukan pembongkaran atau penggantian tanpa inspeksi atau pengukuran pendahulu, misalnya \textit{case} WT34962 dan WT35073. Perbaikan yang relevan adalah mewajibkan setiap pembongkaran, penggantian, atau perbaikan didahului inspeksi atau pengukuran yang terdokumentasi.

Kedua, diagnosis multi-komponen berurutan. Dalam DFG, pada beberapa kesempatan, kategori MEASUREMENT dan INSPECTION disajikan secara berurutan tanpa kategori lain di antaranya. Namun, analisis terhadap atribut komponen dari setiap \textit{event} menunjukkan bahwa pola ini didominasi oleh pemeriksaan atau pengukuran terhadap komponen yang berbeda dalam satu sesi diagnosis: dari enam \textit{event} INSPECTION berturut-turut, semuanya menargetkan komponen yang berbeda, dan dari tujuh \textit{event} MEASUREMENT berturut-turut, lima di antaranya juga menargetkan komponen yang berbeda. Temuan ini menyoroti risiko salah menafsirkan pengulangan kategori tindakan sebagai pengerjaan ulang jika tidak divalidasi pada tingkat komponen, bukan hanya pada tingkat kategori \texttt{action\_type}. Perbaikan yang relevan mempertimbangkan penambahan atribut komponen dalam analisis kinerja DFG sebagai upaya untuk membedakan dengan tepat pola pengerjaan ulang yang sebenarnya (pengulangan pada komponen yang identik) dari diagnosis multi-komponen yang normal dan sehat.

Ketiga, aktivitas setelah pengujian akhir. Pada 4 \textit{case} masih terdapat aktivitas setelah TESTING terakhir, misalnya WT35159 yang diikuti ADJUSTMENT dan WT34939 yang diikuti INSPECTION, yang menandakan pengujian belum lolos. Perbaikan yang relevan adalah mendokumentasikan hasil uji dan akar penyebab kegagalan agar pengujian ulang dapat ditelusuri.

Keempat, jeda panjang sebelum intervensi. Relasi CLEANING $\rightarrow$ ADJUSTMENT dan INSPECTION $\rightarrow$ REPLACEMENT pada relasi lengkap menunjukkan selang penyelesaian yang panjang (59,43 jam dan 38,44 jam), yang mencerminkan waktu tunggu ketersediaan suku cadang atau penjadwalan tim. Perbaikan yang relevan adalah menyiapkan suku cadang kritis lebih awal dan merapatkan penjadwalan menjelang penggantian dan penyetelan.

Keempat pola menunjukkan bahwa \textit{process discovery} tidak sekadar menyederhanakan model, tetapi juga menghasilkan temuan yang dapat ditindaklanjuti untuk memperkuat disiplin urutan kerja, menekan \textit{rework}, dan memperpendek waktu tunggu.
